
\documentclass[../main.tex]{subfiles}
\graphicspath{{\subfix{../images/}}}
\newcommand{\citationneeded}{[\textit{citation needed}]}
% Application of engineering knowledge and understanding and quality of experimental results and/or design work. -- 40% (4,800 words / 5,647.059 words)
%   Expertly applied i) mathematical and statistical methods, tools and notations, ii) quantitative and computational methods, and engineering principles in completing the project. Technical content of the highest quality - could easily provide the basis for a publication or could confidently and proudly be presented to a commercial client. Work of the level of a professional engineer.

%   You should describe how you achieved the project objectives in sufficient detail that the work could be repeated by someone else, perhaps working in a different University.
%   You should include photographs of any equipment used to provide evidence of use.
%   You should present your results in graphical and tabular form providing appropriate numerical and statistical analysis where appropriate.
\usepackage{amssymb}
\usepackage{amsmath}


\begin{document}
\emergencystretch 3em

\section{Methodology and Results}

% \begin{figure*}
%     \centering
%     \includegraphics[width=\linewidth,height=5cm]{example-image-a}
%     \caption{The proposed (a) 1-D CNN, (b) Autoencoder for tool condition monitoring.}
%     \label{fig:cm_models}
% \end{figure*}

This section presents the experimental and theoretical methodologies for data processing in high-precision tool condition monitoring. The theoretical background for two condition monitoring models (1-D CNN and autoencoder) is first described. The models are then validated in software using existing experimental data before being realised in hardware. The hardware implementation is first behaviourally validated, then physically tested.

% The models are illustrated in Figure \ref{fig:cm_models}. Using identical input data, the 1-D CNN model performs tool wear classification, while the autoencoder estimates tool wear.

\subsection{Convolutional Neural Network}

CNNs have gained widespread recognition for their image classification abilities as they are uniquely designed for learning unintuitive spacial relationships in data \citationneeded. Thus, CNNs in various forms have been successfully implemented in condition monitoring applications \cite{RN16}, \cite{RN17}.

The CNN was chosen due to its relative simplicity, prevalence in literature and its inherent efficiency for feature extraction. It serves as a baseline for performance for more or less advanced machine learning techniques.

A typical CNN comprises alternating convolutional and pooling layers, which extract high-level characteristic features from the input data. The specific features extracted depend on the training process and target data, and are numerically determined by the weights of each filter or kernel. For classification tasks, CNNs typically incorporate a Multi-Layer Perceptron (MLP) composed of fully connected layers. At a minimum, one perceptron layer is used as a regression (output) layer. Figure \ref{fig:typical_cnn} shows a typical CNN structure.

\begin{figure}[!ht] 
    \centering 
    \includegraphics[width=0.9\linewidth]{LaTeX/final_report/images/typical_CNN.png} 
    \caption{Typical architecture of a Convolutional Neural Network.} 
    \label{fig:typical_cnn} 
\end{figure}

The difference between a typical CNN architecture and the architecture proposed is the input data format. Rather than using a 2-D image, or time-series data reshaped into a 2-D format, the input array here is a continuous 1-D signal, this type of CNN is known as streaming 1-D convolutional neural network. The advantage of this type of CNN is that, every input sample has a corresponding output, as opposed to every $M$ samples where $M$ is the size of the kernel.

In a CNN, each convolution operation between a kernel and its corresponding signal region is followed by an activation function. The output may be regarded as a series of feature maps. Mathematically, the streaming 1-D convolution layer is described by Equation \ref{eq:conv}:

\begin{equation}
    z_i^l = f(k_i*x+b_i)
    \label{eq:conv}
\end{equation}

where $x\in\mathbb{R}^{1\times L}$ is a one-dimensional input with length $L$, $i=1,2,...,N$ is the index of filters in the convolutional layers, $l$ is the index of the convolutional and pooling layers, $f$ is the activation function, $*$ represents the convolution operation, $b_i$ is the bias belonging to the $i$th filter, $z_i^l$ is the feature generated by the $i$th filter as it travels along the input sequence, $k_i$ is the corresponding kernel matrix of the $i$th filter. The output (feature map) of the $l$th convolutional layer is:

\begin{equation}
    z^l_j=[z^l_1,z^l_2,...,z^l_N]
\end{equation}

where $z^l_j$ is the $j$th feature map of the $l$th layer. Figure \ref{fig:conv} visualises the streaming 1-D convolutional layer.

\begin{figure}[!ht] 
    \centering 
    \includegraphics[width=0.9\linewidth]{LaTeX/final_report/images/conv.png} 
    \caption{Streaming 1-D convolutional layer with one input stream and five filters.} 
    \label{fig:conv} 
\end{figure}

Pooling is commonly used following a convolutional layer. Pooling performs a dimensionality reduction to avoid overfitting and ensure the scale-invariant characteristics of the feature map. The proposed CNN uses global average pooling. Groups of $Q$ features from each filter are averaged, where $Q=2^n$ for some $n\in\mathbb{N}$. Mathematically this is described by Equation \ref{eq:gap}: 

When $j\mod Q=0$, compute:
\begin{equation}
    p^l_i = \frac{1}{Q}\sum_{q=j-Q}^{j}{z^l_{qi}}
    \label{eq:gap}
\end{equation}

where $p^l_i$ is the $i$th pooled feature from the $l$th layer. Using powers of two for pools ensures maximum computational efficiency and, later, eases hardware realisation.

Similar to the convolutional layer, the fully connected layer is described by Equation \ref{eq:fc}:

\begin{equation}
    z^l_i=f(w_i\cdot x+b_i)
    \label{eq:fc}
\end{equation}

where $x\in\mathbb{R}^{1\times L}$ is a one-dimensional input with length $L$, $l$ is the index of the fully connected layer, $f$ is the activation function, $\cdot$ represents the dot product operation, $b_i$ is the bias belonging to the $i$th neuron, $z^l_i$ is the output generated by the $i$th neuron, $w_i$ is the corresponding weight of the $i$th neuron.

% introduce the network and parameters ( do a diagram)

% explain why its small


\subsection{Autoencoder}

\subsection{Training and Testing}

\subsection{Architecture and HW Realisation}

\subsubsection{Neuron Layers}

\subsubsection{ReLU}

\subsubsection{Convolution}

\subsubsection{FFT}

\onecolumn

\biblio

\end{document}
